\documentclass[11pt]{article}
\pagestyle{plain}

% \pgfplotsset{compat=1.18}
\usepackage[margin=1in]{geometry}
\usepackage{amsmath,amsfonts,amssymb,amsthm}
\usepackage[ruled,vlined,linesnumbered,commentsnumbered]{algorithm2e}
\SetKwInput{KwIn}{Input}
\SetKwInput{KwOut}{Output}
\SetKwInput{KwGlobal}{Global variable}
\SetKwInput{KwOracle}{Oracle}
\SetKwComment{Comment}{// }{}
\usepackage{graphicx}
\usepackage{enumerate}
\usepackage{bbm}
\usepackage{verbatim}
\usepackage{hyperref,color}
\usepackage[capitalize,nameinlink]{cleveref}
\usepackage[dvipsnames]{xcolor}
\hypersetup{
	colorlinks=true,
	pdfpagemode=UseNone,
	citecolor=OliveGreen,
	linkcolor=NavyBlue,
	urlcolor=Magenta,
	pdfstartview=FitW
}
\usepackage{appendix}
\usepackage{makecell}
\usepackage{tablefootnote}
\crefname{appsec}{Appendix}{Appendices}
\usepackage{tikz}
\usepackage{pgfplots}
\usepackage{xifthen}
\usepackage{subcaption}
\usepackage{hhline}

\theoremstyle{plain}
\newtheorem{theorem}{Theorem}[section]
\newtheorem{proposition}[theorem]{Proposition}
\newtheorem{lemma}[theorem]{Lemma}
\newtheorem{corollary}[theorem]{Corollary}
\newtheorem{conjecture}[theorem]{Conjecture}
\newtheorem{problem}[theorem]{Problem}
\newtheorem{claim}[theorem]{Claim}
\newtheorem{fact}[theorem]{Fact}

\theoremstyle{definition}
\newtheorem{definition}[theorem]{Definition}
\newtheorem{example}[theorem]{Example}
\newtheorem{observation}[theorem]{Observation}
\newtheorem{assumption}[theorem]{Assumption}
\newtheorem*{assumption*}{Assumption}
\newtheorem{condition}[theorem]{Condition}

\theoremstyle{remark}
\newtheorem{remark}[theorem]{Remark}

\crefname{lemma}{Lemma}{Lemmas}
\crefname{theorem}{Theorem}{Theorems}
\crefname{definition}{Definition}{Definitions}
\crefname{fact}{Fact}{Facts}
\crefname{claim}{Claim}{Claims}
\crefname{proposition}{Proposition}{Propositions}


\newcommand{\dif}{\,\mathrm{d}}
%\newcommand{\E}{\mathbb{E}}
%\newcommand{\Var}{\mathrm{Var}}
\newcommand{\Cov}{\mathrm{Cov}}
\newcommand{\MI}{\mathrm{I}}
%\newcommand{\Ent}{\mathrm{Ent}}
\newcommand{\ind}{\mathbbm{1}}
\newcommand{\allone}{\mathbf{1}}
\newcommand{\tr}{\mathrm{tr}}
\newcommand{\wt}{\mathrm{weight}}
\DeclareMathOperator*{\argmax}{arg\,max}
\newcommand{\norm}[1]{\left\lVert #1 \right\rVert}
\newcommand{\TV}[2]{\left\|#1 - #2\right\|_{\mathrm{TV}}}
\newcommand{\ceil}[1]{\left\lceil #1 \right\rceil}
\newcommand{\floor}[1]{\left\lfloor #1 \right\rfloor}
\newcommand{\ip}[2]{\left\langle #1 , #2 \right\rangle}
\newcommand{\grad}{\nabla}
\newcommand{\hessian}{\nabla^2}
\newcommand{\trans}{\intercal}
\newcommand{\diag}{\mathrm{diag}}
\newcommand{\poly}{\mathrm{poly}}
\newcommand{\dist}{\mathrm{dist}}
\newcommand{\sgn}{\mathrm{sgn}}
\newcommand{\fpras}{\mathsf{FPRAS}}
\newcommand{\fpaus}{\mathsf{FPAUS}}
\newcommand{\fptas}{\mathsf{FPTAS}}

\newcommand{\eps}{\varepsilon}

\newcommand{\N}{\mathbb{N}}
\newcommand{\Z}{\mathbb{Z}}
\newcommand{\Q}{\mathbb{Q}}
\newcommand{\R}{\mathbb{R}}
\newcommand{\C}{\mathbb{C}}

\renewcommand{\AA}{\mathcal{A}}
\newcommand{\BB}{\mathcal{B}}
\newcommand{\CC}{\mathcal{C}}
\newcommand{\DD}{\mathcal{D}}
\newcommand{\EE}{\mathcal{E}}
\newcommand{\FF}{\mathcal{F}}
\newcommand{\GG}{\mathcal{G}}
\newcommand{\HH}{\mathcal{H}}
\newcommand{\II}{\mathcal{I}}
\newcommand{\JJ}{\mathcal{J}}
\newcommand{\KK}{\mathcal{K}}
\newcommand{\LL}{\mathcal{L}}
\newcommand{\MM}{\mathcal{M}}
\newcommand{\NN}{\mathcal{N}}
\newcommand{\OO}{\mathcal{O}}
\newcommand{\PP}{\mathcal{P}}
\newcommand{\QQ}{\mathcal{Q}}
\newcommand{\RR}{\mathcal{R}}
\renewcommand{\SS}{\mathcal{S}}
\newcommand{\TT}{\mathcal{T}}
\newcommand{\UU}{\mathcal{U}}
\newcommand{\VV}{\mathcal{V}}
\newcommand{\WW}{\mathcal{W}}
\newcommand{\XX}{\mathcal{X}}
\newcommand{\YY}{\mathcal{Y}}
\newcommand{\ZZ}{\mathcal{Z}}

\newcommand{\KL}[2]{D_{\mathrm{KL}} \left( #1 \,\Vert\, #2 \right)}
\newcommand{\KLbig}[2]{D_{\mathrm{KL}} \left( #1 \,\big\Vert\, #2 \right)}
\newcommand{\KLBig}[2]{D_{\mathrm{KL}} \left( #1 \,\Big\Vert\, #2 \right)}

\newcommand{\SKL}[2]{D_{\mathrm{SKL}} \left( #1 \,\Vert\, #2 \right)}
\newcommand{\SKLbig}[2]{D_{\mathrm{SKL}} \left( #1 \,\big\Vert\, #2 \right)}
\newcommand{\SKLBig}[2]{D_{\mathrm{SKL}} \left( #1 \,\Big\Vert\, #2 \right)}

\newcommand{\qandq}{\quad\text{and}\quad}
\newcommand{\supp}{\mathrm{supp}}

\newcommand{\DKL}[2]{\-D_{\-{KL}}\left(#1 \parallel #2\right)}
\newcommand{\DTV}[2]{\-D_{\mathrm{TV}}\left({#1},{#2}\right)}
% \newcommand{\dist}{\mathrm{dist}}
\newcommand{\e}{\mathrm{e}}
\renewcommand{\epsilon}{\varepsilon}
\renewcommand{\emptyset}{\varnothing}
% \newcommand{\norm}[1]{\left\Vert#1\right\Vert}
\newcommand{\set}[1]{\left\{#1\right\}}
\newcommand{\tuple}[1]{\left(#1\right)} 
% \newcommand{\eps}{\varepsilon}
\newcommand{\inner}[2]{\left\langle #1,#2\right\rangle}
\newcommand{\tp}{\tuple}
\newcommand{\ol}{\overline}
\newcommand{\abs}[1]{\left\vert#1\right\vert}
\newcommand{\ctp}[1]{\left\lceil#1\right\rceil}
\newcommand{\ftp}[1]{\left\lfloor#1\right\rfloor}
\newcommand{\Ind}[1]{\-{Ind}_{#1}}
\newcommand{\todo}[1]{{\color{red} TODO: #1}}

\newcommand{\Cor}{\Psi^{\mathrm{sym}}}
\newcommand{\cor}{\Psi^{\-{cor}}}
\newcommand{\Inf}{\Psi}
% \newcommand{\diag}{\-{diag}}

\def\*#1{\boldsymbol{#1}} % Use \*A for \mathbf{A}
\def\+#1{\mathcal{#1}} % Use \+A for \mathcal{A}
\def\-#1{\mathrm{#1}} % Use \-A for \mathrm{A}
\def\=#1{\mathbb{#1}} % Use \^A for \mathbb{A}
\def\!#1{\mathfrak{#1}} % Use \!A for \mathfrak{A}

\newcommand*\diff{\mathop{}\!\mathrm{d}}

% probability
% \DeclareMathOperator*{\oPr}{\mathbf{Pr}}
\def\oPr{\mathop{\mathrm{Pr}}}
\renewcommand{\Pr}[2][]{ \ifthenelse{\isempty{#1}}
  {\oPr\left[#2\right]}
  {\oPr_{#1}\left[#2\right]} } % Use \Pr[a]{b} for \mathbf{Pr}_a[b], \Pr{b} for  \mathbf{Pr}[b]

% expectation: relies on the package "dsfont"
% \DeclareMathOperator*{\oE}{\mathbb{E}}
\def\oE{\mathop{\mathbb{E}}}
\newcommand{\E}[2][]{ \ifthenelse{\isempty{#1}}
  {\oE\left[#2\right]}
  {\oE_{#1}\left[#2\right]} }

% variance
%\DeclareMathOperator*{\oVar}{\mathbf{Var}}
\def\oVar{\mathrm{Var}}
\newcommand{\Var}[2][]{ \ifthenelse{\isempty{#1}}
  {\oVar\left[#2\right]}
  {\oVar_{#1}\left[#2\right]} }

% entropy
% \DeclareMathOperator*{\oEnt}{\mathbf{Ent}}
\def\oEnt{\mathrm{Ent}}
\newcommand{\Ent}[2][]{ \ifthenelse{\isempty{#1}}
  {\oEnt\left[#2\right]}
  {\oEnt_{#1}\left[#2\right]} }

\newcommand{\PhiEnt}[2][]{ \ifthenelse{\isempty{#1}}
  {\oEnt^\phi\left[#2\right]}
  {\oEnt^\phi_{#1}\left[#2\right]} }


\newenvironment{Exampleparagraph}[1][]
{
    \par 
    \noindent 
    \textbf{\color{blue} #1}
    \par 
    \vspace{0.5em} 
    \itshape 
}
{
    \par 
    \vspace{1em} 
}

%\input{local-macro}

\newcommand{\RED}[1]{\textcolor{red}{#1}}
\newcommand{\zc}[1]{\textcolor{red}{ZC: #1}}

\newcommand{\bern}[1]{\mathrm{Bern}\left(#1\right)}



\title{Local Bobkov inequality and Eldan-Gross inequality}
\date{\today}

\pgfplotsset{compat=1.18} 
\begin{document}

\maketitle

\section{Local Bobkov inequality}
Let $\mathrm{Sch}(G,X,U)$ be the Schreier graph. We assume that each element in $U$ has exponent $2$, i.e., $g^2 = e$ for all $g \in U$. Furthermore, we assume $g U g^{-1} = U$ for all $g \in U$.
Define the Laplacian operator $L$ as follows:
\begin{align*}
    \forall f:X \to [0,1],\quad Lf(x) = \sum_{g \in U} \tp{f(g x) - f(x)}.
\end{align*}
The Markov semigroup $P_t = \e^{tL}$. We establish the following local Bobkov inequality.
\begin{lemma}[Local Bobkov inequality]
    For any $t \ge 0$ and $f: X \to [0,1]$, it holds
    \begin{align*}
        I(P_t f) \le P_t \tp{\sqrt{I(f)^2 + 2(1-\e^{-4t}) \abs{\nabla f}^2}},
    \end{align*}
    where $\nabla f(x) = (f(g x) - f(x))_{g \in U}$ and $I$ is the Guassian isoperimetric profile.
\end{lemma}
\begin{proof}
    Let $J(s) = P_s F_s$, where $F_s = \sqrt{I(f_{t-s})^2 + 2\tp{1-\e^{-4s}} \abs{\nabla f_{t-s}}^2}$.
    It suffices to prove that $J(s)$ is monotone increasing. 
    Note that
    \begin{align*}
        \frac{\dif}{\dif s} P_s F_s = P_s \tp{(L + \partial_s) F_s}.
    \end{align*}
    Therefore, it suffices to show that $(L+\partial_s) F_s \ge 0$.

    Let $V_s = (I(f_{t-s}), \sqrt{2\tp{1-\e^{-4s}}} \partial_1 f_{t-s}, \ldots, \sqrt{2\tp{1-\e^{-4s}}} \partial_{\abs{U}} f_{t-s})$. 
    It can be seen that $F_s = \abs{V_s}$. By a straightforward calculation, it holds
    \begin{align*}
        L F_s(x) &=  \sum_{g \in U} \tp{|V_s|(g x) - |V_s|(x)}\\
        &= \sum_{g \in U} \tp{\abs{V_s}(g x) - \frac{\abs{V_s}^2(x)}{\abs{V_s}(x)}}\\
        &= \frac{\inner{V_s}{LV_s}}{\abs{V_s}}(x) + \sum_{g \in U} \underbrace{\tp{\abs{V_s}(g x) - \frac{\inner{V_s(x)}{V_s(g x)} }{\abs{V_s}(x)}}}_{:=D_g}.
    \end{align*}
    Furthermore,
\begin{align*}
\partial_s F_s(x) = \partial_s \sqrt{\sum_{i=0}^n V_{s,i}^2} = \frac{\inner{V_{s}}{\partial_s V_s}}{\abs{V_s}}(x). 
\end{align*}
Hence, we have
\begin{align*}
    (\partial_s + L) F_s(x) = \frac{\inner{V_s}{(\partial_s + L)V_s}}{\abs{V_s}} + \sum_{g \in U} D_g.
\end{align*}
By a straightforward calculation, we have
\begin{align*}
    (\partial_s+L) \sqrt{2(1-\e^{-4s})} \partial_i f_{t-s} = \frac{4\e^{-4s}}{\sqrt{2(1-\e^{-4s})}} \partial_i f_{t-s},
\end{align*}
where we use the fact that $g U g^{-1} = U$ for all $g \in U$.
Furthermore, note that
\begin{align*}
    (\partial_{g} f_{t-s})(gx) = f_{t-s}(x) - f_{t-s}(g x) = - \partial_g f_{t-s}(x).
\end{align*}
Let $\widetilde{V}_s(x)$ be the vector $V_s(x)$ except $(\widetilde{V}_s(x))_g = - (V_s(x))_g$.
Hence, the quantity $D_g$ satisfies
\begin{align*}
    D_g &= \abs{V_s}(gx) - \frac{\inner{V_s(x)}{V_s(gx)}}{\abs{V_s}(x)}\\
    &=\abs{V_s}(gx) - \frac{\inner{\widetilde{V}_s(x)}{V_s(gx)}}{\abs{\widetilde{V}_s}(x)} + \frac{2(V^2_s(x))_g}{\abs{V_s}(x)}\\
    &\ge  \frac{2(V^2_s(x))_g}{\abs{V_s}(x)}
\end{align*}
Therefore, it holds
\begin{align*}
    (\partial_s + L) F_s(x) &=\frac{1}{\abs{V_s}}\tp{I(f_{t-s}) \cdot (\partial_s + L) I(f_{t-s}) + 4 \e^{-2s} \abs{\nabla f_{t-s}}^2} + \sum_{g \in U} D_g\\
    &\ge \frac{1}{\abs{V_s}}\tp{I(f_{t-s}) \cdot (\partial_s + L) I(f_{t-s}) + 4 \abs{\nabla f_{t-s}}^2}.
\end{align*}
Finally, we have
\begin{align*}
    I(f_{t-s}) (\partial_s + L) I(f_{t-s}) = I(f_{t-s}) \sum_{g \in U}\tp{I(f_{t-s}(gx)) - I(f_{t-s}) -I'(f_{t-s}) \partial_g f_{t-s}} \ge - \abs{\nabla f_{t-s}}^2,
\end{align*}
where we use the inequality $I(a) \tp{I(b) - I(a) - I'(a)(b-a)} \ge -(b-a)^2$.
\end{proof}

\begin{remark}
    Let $X = \binom{[n]}{k}$ and $U$ be the collection of transpositions. 
    We have $g^2 = id$ for all $g \in U$ and $g U g^{-1} = U$.
\end{remark}
\section{Eldan--Gross inequality}
In this section, we prove the Eldan-Gross inequality for slices.
Let $X = \binom{[n]}{k}$ and $U$ be the collection of transpositions. 
Let $N = \abs{U} = \binom{n}{2}$ and $\rho = \frac{1}{n \log (1/\nu(k))}$ where $\nu(k) = \frac{k(n-k)}{N}$.
\begin{theorem}[Eldan-Gross inequality]\label{thm:EG-ineq}
    For any $f:X \to \{-1,+1\}$, it holds
    \begin{align*}
        \E[]{\abs{\nabla f}} \gtrsim \sqrt{N \rho} \Var[\mu]{f} \sqrt{\log\tp{1+ \frac{\e}{W(f)}}},
    \end{align*}
    where $W(f) = \sum_{\tau \in U} I^2_\tau(f)$ and $I_\tau(f) = \E[\mu]{\abs{\partial_\tau f}}$.
    Here, $\partial_\tau f(x) = \tp{f(\tau x) - f(x)}/2$.
\end{theorem}

\paragraph{Step 1. Local Bobkov inequality to Improved Talagrand's inequality}
\begin{lemma}[Improved Talagrand's inequality]
    For any function $f : X \to \{-1,+1\}$, it holds
    \begin{align}\label{eq:improved-talagrand}
        \E[]{\abs{\nabla{f}}} \gtrsim \sqrt{N \rho} \cdot \Var[]{f} \sqrt{\log \frac{\e}{\Var[]{f}}}.
    \end{align}
\end{lemma}
\begin{proof}
    Let $h = \frac{1+f}{2}$ and $H_t = P_{t / N}$. By local Bobkov inequality, it holds
\begin{align*}
    I\tp{\frac{1+H_t f}{2}} \le \sqrt{2 \tp{1-\e^{-4t/N}}} H_t \abs{\nabla{f}}.
\end{align*}
Together with $I\tp{\frac{1+s}{2}} \gtrsim (1-s^2)$, it holds
\begin{align*}
    1-(H_t f)^2 \lesssim \sqrt{1-\e^{-4t/N}} H_t \abs{\nabla f}.
\end{align*}
Taking expectation over $\mu$ on both sides, we have
\begin{align}\label{eq:var-contraction}
    \Var[]{f} - \Var[]{H_t f} \lesssim \sqrt{1-\e^{-4t/N}} \E[]{\abs{\nabla f}}.
\end{align} 
We first assume $\Var[]{f} \ge \frac{1}{\e^4}$. Plugging $t = 1/\rho$, the variance $\Var[]{H_{1/\rho}}$ satisfies
\begin{align*}
    \Var[]{H_{1/\rho} f} \le \e^{-2} \Var[]{f}
\end{align*}
Combining with~\eqref{eq:var-contraction}, it holds
\begin{align*}
    \Var[]{f} \lesssim \sqrt{1-\e^{-4/\rho N}} \E[]{\abs{\nabla f}} \lesssim \frac{1}{\sqrt{\rho N}} \E[]{\abs{\nabla{f}}}.
\end{align*} 
Together with the lower bound $\Var[]{f} \ge \frac{1}{\e^4}$, we establish~\eqref{eq:improved-talagrand}. For the case where $\Var[]{f} \le \frac{1}{\e^4}$, we use the hypercontractivity to bound $\Var[]{H_t f}$.
Let $g = f - \E[]{f}$. 
By hypercontractivity, we have
\begin{align*}
    \Var[]{H_t f} = \Var[]{H_t g} = \norm{H_t g}_2^2 \le \norm{g}_{p_t}^2 \le \norm{g}_1^{2\theta_t} \norm{g}_2^{2-2\theta_t} = \tp{\Var[\mu]{f}}^{1+\theta_t},
\end{align*}
where $p_t = 1+\e^{-2\rho t}$ and $\theta_t = \frac{1-\e^{-2\rho t}}{1+\e^{-2\rho t}}$.
Therefore, we have
\begin{align}\label{eq:contract}
    \Var[]{f} - \tp{\Var[]{f}}^{1+\theta_t} \lesssim \sqrt{1-\e^{-4t/N}} \E[]{\abs{\nabla f}}.
\end{align}
We now fix $\eps = \frac{1}{\log \tp{\e / \Var[]{f}}} \le \frac{1}{3}$, and $1-\e^{-2\rho t} = \eps$. In this case, $\tp{\Var[]{f}}^{\theta_t} \le \tp{\Var[]{f}}^{\epsilon/2} \le 1/\e$.
Therefore, we have
\begin{align*}
   \Var[]{f} \lesssim \sqrt{1-\e^{-4t/N}} \E[]{\abs{\nabla f}}.
\end{align*}
On the other hand,
\begin{align*}
    1-\e^{-4t/N} \lesssim \frac{4\epsilon}{N \rho}.
\end{align*}
Therefore, we have
\begin{align*}
    \E[]{\abs{\nabla f}} \gtrsim \sqrt{N \rho} \cdot \Var[]{f} \sqrt{\log \tp{\frac{\e}{\Var[]{f}}}}.
\end{align*}
\end{proof}

\paragraph{Step 2. Smoothing lemma}
\begin{lemma}\label{lem:smooth}
    For any function $f: X \to \{-1,+1\}$, it holds
    \begin{align*}
        \Var[]{H_t f} \le W(f)^{\theta_t} \Var[]{f}^{1-\theta_t},
    \end{align*}
    where $\theta_t = \frac{1-\e^{-2\rho t}}{1+\e^{-2\rho t}}$ and $W(f)$ is the sum of squared influences.
\end{lemma}
\begin{proof}
    Let $S_n^{(m)}$ denote the set of permutations $\tau \in S_n$ such that for all $i \le m$, it holds $\tau_i = i$. Let $Q_m$ be the averaging operator defined as 
    \begin{align*}
        (Q_m f)(x) = \frac{1}{(n-m)!}\sum_{\tau \in S_n^{(m)}} f(\tau x).
    \end{align*}
    Specifically, $(Q_m f)(x)$ is the average of values $f(y)$ over $y$'s satisfying $i \in x$ if and only if $i \in y$ for all $i \le m$.
    Fix $f: X \to \{-1,+1\}$, the variance of $H_t f$ satisfies
    \begin{align*}
        \Var[]{H_t f} = \norm{\sum_{i=1}^n \tp{Q_i H_t f - Q_{i-1} H_t f}}_2^{2} &= \sum_{i=1}^n \norm{Q_i H_t f - Q_{i-1} H_t f}_2^2\\
        (\text{$L Q_i = Q_i L$})\quad &=\sum_{i=1}^n \norm{H_t (Q_i - Q_{i-1}) f}_2^2\\
        (\text{hypercontractivity}) \quad &\le \sum_{i=1}^n \norm{(Q_i - Q_{i-1}) f}^2_{q_t}\\
        &\le \sum_{i=1}^n \norm{(Q_i - Q_{i-1}) f}_{1}^{2 \theta_t} \norm{(Q_i - Q_{i-1}) f}_{2}^{2- 2 \theta_t}\\
        &\le \tp{\sum_{i=1}^n \norm{(Q_i-Q_{i-1}) f}_1^2}^{\theta_t} \cdot \Var[]{f}^{1-\theta_t}.
    \end{align*} 
    Note that we are able to generate a uniform sample $\tau \in S_n^{(i-1)}$ by first generating $\tau \in S_n^{(i)}$ uniformly at random, choosing an index $j \ge i$ independently at random, and finally swapping $\tau_i$ and $\tau_j$.
    Therefore, the quantity $\norm{(Q_i - Q_{i-1}) f}_1$ satisfies
    \begin{align*}
        \norm{(Q_i - Q_{i-1}) f}_1 \le \frac{1}{n-i+1} \sum_{\tau = (i,j) \in U}\norm{\abs{\partial_\tau f}}_1 =  \frac{1}{n-i+1} \sum_{\tau = (i,j) \in U} I_\tau(f).
    \end{align*}
    Therefore, we have the following:
    \begin{align*}
        \norm{(Q_i-Q_{i-1}) f}_1^2 \le \frac{(n-i)}{(n-i+1)^2} \sum_{\tau = (i,j) \in U} I^2_\tau(f).
    \end{align*}
    This concludes the proof.
\end{proof}

\paragraph{Step 3. Wrapping up}
We now ready to prove~\Cref{thm:EG-ineq}. 
It suffices to show~\Cref{thm:EG-ineq} when $W(f) < \frac{1}{100}$ and $\Var[]{f} \ge \max\set{\sqrt{W(f)}, 100W(f)}$.
It then follows from~\eqref{eq:contract},~\Cref{lem:smooth} and the choice $\eps = \frac{1}{\log \tp{\e/W(f)}}$.

\section{Eldan--Gross Implies KKL}
Suppose $f: \{0,1\}^n \to \{+1, -1\}$.

Definitions:
\begin{align*}
    V &= \Var[]{f}\\
    I &= \sum_{i=1}^n \mathrm{Inf}_i(f)\\
    W &= \sum_{i=1}^n \mathrm{Inf}_i(f)^2
\end{align*}

\begin{lemma}
    Suppose $f: \{0,1\}^n \to \{+1, -1\}$.
    \begin{align}
        &\text{Eldan--Gross inequality} \nonumber\\
        \quad\Longrightarrow\quad
        & V^2 \log\left(1 + \frac{e}{W}\right) \lesssim I \label{eq:EG-coro}\\
        \quad\Longrightarrow\quad
        & V \log\left(1 + \frac{eV}{W}\right) \lesssim I 
    \end{align}
\end{lemma}

\begin{proof}
    Take 
    \begin{align*}
        \ell \approx \frac{1}{V_f}.
    \end{align*}
    Define $g: \{0,1\}^{n\ell} \to \{+1, -1\}$, such that for every $x = (x^{(1)}, x^{(2)}, \dots, x^{(\ell)}) \in \{0,1\}^{n\ell}$ where each $x^{(i)} \in \{0,1\}^n$, 
    \begin{align*}
        g(x^{(1)}, x^{(2)}, \dots, x^{(\ell)}) = \prod_{i=1}^\ell f(x^{(i)}).
    \end{align*}

    We observe that
    \begin{align*}
        V_g &= 1 - (1 - V_f)^\ell = \Theta(1); \\
        I_g &= \ell I_f; \\
        W_g &= \ell W_f.
    \end{align*}
    For the variance, this follows from that $V_g = 1 - (\mathbb{E}[g])^2 = 1 - (\mathbb{E}[f])^{2\ell} = 1 - (1 - V_f)^\ell$.
    \begin{align*}
        &V_g^2 \log\left(1 + \frac{e}{W_g}\right) \lesssim I_g \tag{\cref{eq:EG-coro} for $g$}\\
        \Longrightarrow \qquad
        &\log\left(1 + \frac{e V_f}{W_f}\right) \lesssim \frac{I_f}{V_f} 
        %\Longrightarrow \qquad
        %&V_f^2 \log^2\left(1 + \frac{e V_f}{W_f}\right) \lesssim n W_f \tag{$I^2 \le n W$}
    \end{align*}
\end{proof}

\begin{lemma}
    Suppose $f: \{0,1\}^n \to \{+1, -1\}$.
    Let $M = \max_i \mathrm{Inf}_i(f)$.
    \begin{align}
        V \log\left(1 + \frac{eV}{W}\right) \lesssim I 
        \quad\Longrightarrow\quad
        \begin{cases}
            \frac{\log^2 n}{n} V^2 \lesssim W \\
            V \log\left(1 + \frac{1}{M}\right) \lesssim I 
        \end{cases} 
        \quad\Longrightarrow\quad
        M \gtrsim \frac{\log n}{n} V \tag{KKL}
    \end{align}
\end{lemma}

\section*{Questions}
\begin{itemize}
\item Is there any application of Eldan-Gross inequality?
\item I suspect this approach is insufficient to establish a local Bobkov inequality for mixing Gibbs measures on $\{-1,+1\}^n$. Although I don’t have immediate evidence to confirm this.
\end{itemize}
\end{document}

